% Minimal retained preamble for lane-p2 cards taken from
% sources/logic-fl02-conservation-of-ramsey-s-theorem-f/pinned-original-source.tex
% (Conservation of Ramsey's theorem for pairs, arXiv:2402.11616v3).
%
% Copied verbatim from the author's own file:
%   line 2      : \documentclass{article}
%   lines 16-25 : the author's theorem declarations
%   lines 46-88 : the author's macro definitions for the reverse-mathematics principles
%   line 91     : \card
%   lines 104-107: the natural sums and products, which rely on mathabx (source line 103)
%   line 116    : \qt
%
% Omitted author preamble lines and why (nothing was renamed or redefined):
%   line 3 (graphicx), 7 (url), 9 (cancel): no figure, url or cancel use in the retained spans;
%   line 6 (cleveref): the assembler adds its own cleveref machinery, a second load clashes;
%   line 8 (\usepackage[utf8x]{inputenc}): inputenc with the utf8x option fights the forced
%       xelatex engine;
%   lines 28-37 (the rep@theorem machinery and \newreptheorem{maintheorem}): the retained spans
%       use no rep-theorem environment;
%   lines 40-43 (xcolor and the three author comment macros): not used by the retained spans;
%   lines 93-99 (the LS2 font encoding block for \smashtimes based on stix): the host TeX Live has
%       no stix directory and the author himself never uses the result -- line 105 maps \ntimes to
%       \dottimes, and line 110 shows \smashtimes already commented out in the source.
% mathabx is retained because \dotplus, \dottimes and \dotdiv come from it.
\documentclass{article}
\usepackage{amsmath,amssymb,amsthm}
\usepackage{mathabx}

\newtheorem{theorem}{Theorem}
\numberwithin{theorem}{section}
\newtheorem{maintheorem}[theorem]{Main Theorem}
\newtheorem{lemma}[theorem]{Lemma}
\newtheorem{question}[theorem]{Question}
\newtheorem{proposition}[theorem]{Proposition}
\newtheorem{statement}[theorem]{Statement}
\newtheorem{remark}[theorem]{Remark}
\newtheorem{definition}[theorem]{Definition}
\newtheorem{corollary}[theorem]{Corollary}

\newcommand{\RCA}{\mathsf{RCA}}
\newcommand{\WKL}{\mathsf{WKL}}
\newcommand{\SRT}{\mathsf{SRT}}
\newcommand{\EM}{\mathsf{EM}}
\newcommand{\SEM}{\mathsf{SEM}}
\newcommand{\ADS}{\mathsf{ADS}}
\newcommand{\SCAC}{\mathsf{SCAC}}
\newcommand{\SADS}{\mathsf{SADS}}
\newcommand{\COH}{\mathsf{COH}}
\newcommand{\ACA}{\mathsf{ACA}}
\newcommand{\qvdash}{\operatorname{{?}{\vdash}}}
\newcommand{\nqvdash}{\operatorname{{?}{\nvdash}}}
\newcommand{\BSig}{\mathsf{B}\Sigma^0}
\newcommand{\ISig}{\mathsf{I}\Sigma^0}
\newcommand{\LSig}{\mathsf{L}\Sigma^0}
\newcommand{\BPi}{\mathsf{B}\Pi^0}
\newcommand{\IPi}{\mathsf{I}\Pi^0}
\newcommand{\LPi}{\mathsf{L}\Pi^0}
\newcommand{\IDel}{\mathsf{I}\Delta^0}
\newcommand{\BCDe}{\mathsf{BC}\Delta^0}
\newcommand{\WF}{\mathsf{WF}}
\newcommand{\RT}{\mathsf{RT}}
\newcommand{\BME}{\mathsf{BME}}
\newcommand{\TT}{\mathsf{TT}}
\renewcommand{\O}{\mathcal{O}}
\renewcommand{\L}{\mathcal{L}}
\newcommand{\C}{\mathcal{C}}
\newcommand{\R}{\mathcal{R}}
\renewcommand{\S}{\mathcal{S}}
\newcommand{\T}{\mathcal{T}}
\renewcommand{\P}{\mathcal{P}}
\newcommand{\M}{\mathcal{M}}
\newcommand{\Nc}{\mathcal{N}}
\renewcommand{\P}{\mathcal{P}}
\newcommand{\Q}{\mathcal{Q}}
\newcommand{\NN}{\mathbb{N}}
\newcommand{\PP}{\mathbb{P}}
\newcommand{\uh}[0]{{\upharpoonright}}
\newcommand{\card}{\operatorname{card}}
\newcommand{\nplus}{\dotplus}
\newcommand{\ntimes}{\dottimes}
\newcommand{\nminus}{\dotdiv}
\newcommand{\nsum}{\dot{\sum}}
\def\qt#1{``#1''}
\title{Shore blocking turns one descending sequence below an ordinal into one below its natural product}
\author{Quentin Le Houérou, Ludovic Levy Patey, Keita Yokoyama}
\date{arxiv:2402.11616v3; CC BY 4.0}
\begin{document}
\maketitle
\noindent\textit{Editorial note.} The counted unit is the article's own Lemma lem:shore-blocking-wf at source lines 426-431 together with the proof environment that begins two lines after the statement ends (433-450); the pairing is the ordinary adjacent one.\par\medskip
\noindent\textit{Editorial note.} The retained context at source line 390 is the author's own announcement of the two combinatorial lemmas and of the natural (Hessenberg) operations; the retained context at source lines 392-398 is his own definition of the natural sum and of the natural product alpha ntimes k.\par\medskip
\noindent\textit{Editorial note.} The retained preamble reproduces verbatim the author's document class at source line 2, his theorem declarations at source lines 16-25, his macro definitions at source lines 46-88 and 91, his natural-operation macros at source lines 104-107 and his quotation macro at source line 116. The LS2 font-encoding block at source lines 93-99, which declares \textbackslash{}smashtimes through the stix fonts, is not retained because the host TeX Live has no stix directory and because the author never uses that symbol: source line 105 defines \textbackslash{}ntimes as \textbackslash{}dottimes and source line 110 shows the alternative already commented out. mathabx (source line 103) is retained because \textbackslash{}dotplus, \textbackslash{}dottimes and \textbackslash{}dotdiv come from it. The cleveref load at source line 6 is not retained because the assembler adds its own cleveref machinery, and the inputenc load with the utf8x option at source line 8 is not retained because it fights the forced xelatex engine. Nothing was renamed, redefined or added.\par\medskip
\noindent\textit{Editorial note.} No retained line contains a citation or a cross-reference, so the material carries no bibliography and needs no navigation pointer. In particular the sentence of the introduction that would name the proof-theoretic ordinal references is not part of the retained context.\par\medskip
\noindent\textit{Editorial note.} The article's conservation theorems, its other combinatorial lemma lem:wf-product, its lemma lem:coh-wf-forcing-jump and lem:coh-wf-forcing-wf, and its proposition prop:coh-wf are not counted here and are not used as premises.\par\medskip
\noindent\textit{Allowed standard prerequisites (editorial).}\begin{itemize}
\item Natural sum and natural product of ordinals below epsilon\_0. For ordinals alpha and beta below epsilon\_0 written with a common family of exponents, the natural sum is obtained by adding the coefficients of like powers of omega, and alpha ntimes k is the natural sum of alpha with itself k times; the article's own definition is retained at source lines 392-398.. Fixes the two operations that the counted lemma uses in both its hypothesis and its conclusion, namely beta below alpha and the larger ordinal beta ntimes k.
\item Theories RCA\_0, I Sigma\_1 and the well-foundedness statement WF(alpha). RCA\_0 is the standard base theory of reverse mathematics; I Sigma\_1 is the induction scheme for Sigma\textasciicircum{}0\_1 formulas; WF(alpha) states that the ordinal alpha is well founded, and a model M satisfies WF(alpha) when it has no M-infinite descending sequence of ordinals below alpha; these are the standard notions of the reverse-mathematics literature and the article uses the notations RCA\_0, I Sigma\_1 and WF(alpha) throughout, for example in its abstract and introduction.. The counted lemma is stated for a model M of RCA\_0 plus WF(alpha) and its proof invokes I Sigma\_1 to make the constructed Turing functional total, so both theories and the meaning of the well-foundedness statement must be available.
\item M[X] and the forcing extension by a real. For a model M = (M,S) of second-order arithmetic and a real X in 2\textasciicircum{}M, the article writes M[X] for the extension of M by the real X; the counted lemma is stated for every X with M[X] a model of RCA\_0.. The hypothesis and the proof both quantify over reals X such that the extension M[X] satisfies RCA\_0, and the proof appeals to M[X] satisfying I Sigma\_1.
\end{itemize}\medskip
Before working on $\Pi^1_1$-conservation of~$\WKL$ and $\COH$, we prove two combinatorial lemma which say that Shore blocking comes for free for preserving $\WF(\alpha)$ for any ordinal $\alpha \leq \epsilon_0$. Instead of using the standard sums and products over ordinals, the natural (Hessenberg) sums and products will be more convenient:

\begin{definition}[Natural sum and product]
Let $\alpha$ and $\beta$ be two ordinals less than $\epsilon_0$. Let $\alpha = \omega^{\gamma_1}n_1 + \dots + \omega^{\gamma_k}n_k$ and $\beta = \omega^{\gamma_1}m_1 + \dots + \omega^{\gamma_k}m_k$ (We allow the $n_i$ and $m_i$ to be equal to $0$ in order to write $\alpha$ and $\beta$ using the same exponents $\gamma_i$). We define the natural sum of $\alpha$ and $\beta$ (denoted by $\alpha \nplus \beta$) to be the following ordinal (It does not depend on the choice of the family $(\gamma_i)$) : 

$$\omega^{\gamma_1}(n_1 + m_1) + \dots + \omega^{\gamma_k}(n_k + m_k)$$

We also define $\alpha \ntimes k$ to be equal to be the natural sum of $\alpha$ with itself $k$ times. %and $\alpha \ntimes \omega = \omega^{\gamma_1+1}n_1 + \dots + \omega^{\gamma_k+1}n_k$ (whereas $\alpha \times \omega = \omega^{\gamma_1 + 1}n_1$)
\end{definition}

\begin{lemma}\label[lemma]{lem:shore-blocking-wf}
Consider a model $\M = (M,S)$ and $\alpha \in M$ an ordinal less than $\epsilon_0$.
For every~$k \in M$, there is a Turing functional $\Gamma_k$ such that, letting~$\beta \in \alpha$ be the largest ordinal less than $\alpha$ of index at most~$k$, for every~$X \in 2^M$ such that $\M[X] \models \RCA_0$, if there is some~$e < k$ such that $\Phi^X_e$ is an $M$-infinite decreasing sequence of elements smaller than~$\beta$, then $\Gamma^X_k$ is an $M$-infinite decreasing sequence of elements smaller than~$\beta \ntimes k$.

Moreover, an index of~$\Gamma_k$ can be found computably in~$k$.
\end{lemma}

\begin{proof}
By twisting the Turing functionals, we can assume that for every~$e, n \in M$, if $\Phi^\sigma_e(n)\downarrow$, then
\begin{itemize}
    \item[(1)] $n < |\sigma|$ ;
    \item[(2)] $\Phi^\sigma_e(m)\downarrow$ for every~$m < n$ ;
    \item[(3)] $\Phi^\sigma_e(0), \Phi^\sigma_e(1), \dots, \Phi^\sigma_e(n)$ is a strictly decreasing sequence of elements smaller than~$\beta$.
\end{itemize}

Given $\sigma \in 2^{<M}$ and $e < k$, let $\alpha(\sigma, e) = \Phi^{\sigma}_e(s)$ for the largest $s < |\sigma|$ such that $\Phi^{\sigma}_e(s)\downarrow$. If there is no such~$s$, then $\alpha(\sigma, e) = \beta$. Note that if $\sigma' \succeq \sigma$, then $\alpha(\sigma',e) \leq \alpha(\sigma,e)$.

Let~$\sigma_{-1} = \epsilon$.
Let~$\Gamma^X_k$ be the Turing functional which, on input~$n$, searches for some~$x > |\sigma_{n-1}|$ and some~$\sigma_n \prec X$ such that $\Phi_e^{\sigma_n}(x)\downarrow$ for some~$e < k$. If found, it outputs $\alpha(\sigma, 0) \nplus \dots \nplus \alpha(\sigma, k-1)$. Note that if~$\Gamma^X_k(n)\downarrow$, then by (3), $\Gamma^X(n)$ is an ordinal smaller than $\beta \ntimes k$.

Suppose that $X$ is such that $\M[X] \models \RCA_0$ and there is an $e < k$ is such that $\Phi_e^X$ is total.
Since~$\M[X] \models \ISig_1$, $\Gamma^X$ is total.

Moreover, since~$x > |\sigma_{n-1}|$, letting~$e < k$ be such that $\Phi_e^{\sigma_n}(x)\downarrow$, by (1) we have $\Phi_e^{\sigma_{n-1}}(x)\uparrow$, so by (2) and (3), $\alpha(\sigma_{n+1},e) < \alpha(\sigma_n, e)$, hence $\Gamma^X_k(n+1) < \Gamma^X_k(n)$. It follows that $\Gamma^X_k$ is an $M$-infinite decreasing sequence of ordinals smaller than~$\beta \ntimes k$.
\end{proof}

\end{document}
